\section{Estado del Arte}
La previsión financiera personal a través de series temporales conforma en la actualidad un ecosistema dinámico investigado activamente. Numerosas propuestas proponen aplicaciones donde, dada una serie de movimientos, algoritmos de \textit{Machine Learning} son desplegados para detectar el perfil del comprador \cite{online_shopper} o pronosticar la evolución \cite{personal_finance_ml}.

Históricamente, los enfoques algorítmicos recayeron fuertemente en aproximaciones estadístico-agresoras como los modelos ARIMA \cite{box2015time}. Este tipo de métodos sigue vigente hoy dada su robustez frente al error a largo plazo. En los desarrollos modernos paralelos para aplicaciones de banca e inteligencia de negocios como \textit{WONGA} \cite{wonga2023}, las aproximaciones ensamble tipo Gradient Boosting (como XGBoost o HistGradientBoosting) han despuntado ofreciendo una notable flexibilización sobre variables no lineales del ámbito temporal (cambios de hábito de la persona) \cite{personal_finance_tracker}.
